%% ============================================================================
%% preamble_shared.tex
%% Shared preamble of main_en.tex and main_zh.tex. Everything that must be identical across the two language
%% versions lives here, so that section files, labels, theorem counters, number macros and notation stay
%% mirrored by construction. main_zh.tex defines the Chinese head names before it reads this file.
%% ============================================================================

%% ---- Latin text font -------------------------------------------------------
%% Under pdflatex, main_en.tex loads the times package. Under xelatex, TeX Gyre Termes supplies the Times
%% metrics when it is installed. fontspec is never loaded under pdflatex.
\usepackage{iftex}
\ifPDFTeX
  \usepackage[T1]{fontenc}
\else
  \usepackage{fontspec}
  \IfFontExistsTF{texgyretermes-regular.otf}{%
    \setmainfont{texgyretermes}[Extension=.otf, UprightFont=*-regular, BoldFont=*-bold,
      ItalicFont=*-italic, BoldItalicFont=*-bolditalic]%
  }{}
\fi

%% ---- packages --------------------------------------------------------------
\usepackage{amsmath,amssymb,amsthm}
\usepackage{booktabs,array,adjustbox,multirow,colortbl,pifont}
\usepackage{graphicx}
\usepackage{xcolor}
\usepackage{algorithm}
\usepackage{algorithmic}
\renewcommand{\algorithmicrequire}{\textbf{Input:}}
\renewcommand{\algorithmicensure}{\textbf{Output:}}
\usepackage{subcaption}
\usepackage{url}
\usepackage{hyperref}
\definecolor{citecolor}{HTML}{009B3A}
\definecolor{refcolor}{HTML}{4527C4}
\hypersetup{colorlinks=true, citecolor=citecolor, linkcolor=refcolor, urlcolor=refcolor, pdfborder={0 0 0}}

%% ---- table style (canonical copy in figure-style/tables/table_macros.tex) --
% Table style of figure-style/tables, one palette shared by the main, capability and dataset tables.
% Put % Table style of figure-style/tables, one palette shared by the main, capability and dataset tables.
% Put % Table style of figure-style/tables, one palette shared by the main, capability and dataset tables.
% Put \input{table_macros} in the preamble. Needs booktabs, colortbl, xcolor, array, pifont.
%
%   first place in a column   \best{0.957}         dark background, bold
%   second place in a column  \secondbest{0.847}   light background
%   family row of main table  \grp{\ncol}{Deep matching} \\
%   row of the proposed method \oursrow before the row
%   capability marks          \cmark  \xmark
%   value waiting for backfill \pend{0.xx}          red, must be gone before submission

\definecolor{tabBest}{RGB}{31,78,121}     % first place, used at 35 percent tint
\definecolor{tabSecond}{RGB}{207,224,237} % second place
\definecolor{tabBand}{RGB}{226,234,242}   % row of the proposed method

\newcommand{\best}[1]{\cellcolor{tabBest!35}\textbf{#1}}
\newcommand{\secondbest}[1]{\cellcolor{tabSecond}#1}
\newcommand{\grp}[2]{\multicolumn{#1}{l}{\textbf{\textit{#2}}}}
\newcommand{\oursrow}{\rowcolor{tabBand}}
\newcommand{\cmark}{\ding{51}}
\newcommand{\xmark}{\ding{55}}
\newcommand{\pend}[1]{\textcolor{red}{#1}}
 in the preamble. Needs booktabs, colortbl, xcolor, array, pifont.
%
%   first place in a column   \best{0.957}         dark background, bold
%   second place in a column  \secondbest{0.847}   light background
%   family row of main table  \grp{\ncol}{Deep matching} \\
%   row of the proposed method \oursrow before the row
%   capability marks          \cmark  \xmark
%   value waiting for backfill \pend{0.xx}          red, must be gone before submission

\definecolor{tabBest}{RGB}{31,78,121}     % first place, used at 35 percent tint
\definecolor{tabSecond}{RGB}{207,224,237} % second place
\definecolor{tabBand}{RGB}{226,234,242}   % row of the proposed method

\newcommand{\best}[1]{\cellcolor{tabBest!35}\textbf{#1}}
\newcommand{\secondbest}[1]{\cellcolor{tabSecond}#1}
\newcommand{\grp}[2]{\multicolumn{#1}{l}{\textbf{\textit{#2}}}}
\newcommand{\oursrow}{\rowcolor{tabBand}}
\newcommand{\cmark}{\ding{51}}
\newcommand{\xmark}{\ding{55}}
\newcommand{\pend}[1]{\textcolor{red}{#1}}
 in the preamble. Needs booktabs, colortbl, xcolor, array, pifont.
%
%   first place in a column   \best{0.957}         dark background, bold
%   second place in a column  \secondbest{0.847}   light background
%   family row of main table  \grp{\ncol}{Deep matching} \\
%   row of the proposed method \oursrow before the row
%   capability marks          \cmark  \xmark
%   value waiting for backfill \pend{0.xx}          red, must be gone before submission

\definecolor{tabBest}{RGB}{31,78,121}     % first place, used at 35 percent tint
\definecolor{tabSecond}{RGB}{207,224,237} % second place
\definecolor{tabBand}{RGB}{226,234,242}   % row of the proposed method

\newcommand{\best}[1]{\cellcolor{tabBest!35}\textbf{#1}}
\newcommand{\secondbest}[1]{\cellcolor{tabSecond}#1}
\newcommand{\grp}[2]{\multicolumn{#1}{l}{\textbf{\textit{#2}}}}
\newcommand{\oursrow}{\rowcolor{tabBand}}
\newcommand{\cmark}{\ding{51}}
\newcommand{\xmark}{\ding{55}}
\newcommand{\pend}[1]{\textcolor{red}{#1}}

\makeatletter\let\inputrows\@@input\makeatother

%% ---- theorem-like environments ---------------------------------------------
%% Head names go through macros, so both language versions share counters and numbering.
\providecommand{\headtheorem}{Theorem}
\providecommand{\headlemma}{Lemma}
\providecommand{\headproposition}{Proposition}
\providecommand{\headcorollary}{Corollary}
\providecommand{\headdefinition}{Definition}
\providecommand{\headassumption}{Assumption}
\providecommand{\headexample}{Example}
\providecommand{\headremark}{Remark}
\providecommand{\headfinding}{Finding}
\providecommand{\headproof}{Proof}

\theoremstyle{plain}
\newtheorem{theorem}{\headtheorem}[section]
\newtheorem{lemma}{\headlemma}[section]
\newtheorem{proposition}{\headproposition}[section]
\newtheorem{corollary}{\headcorollary}[section]
\theoremstyle{definition}
\newtheorem{definition}{\headdefinition}[section]
\newtheorem{assumption}{\headassumption}[section]
\newtheorem{example}{\headexample}[section]
%% Findings are numbered through the whole paper (Finding 1 to Finding 7), not per section.
\newtheorem{finding}{\headfinding}
\theoremstyle{remark}
\newtheorem{remark}{\headremark}[section]
\renewcommand{\proofname}{\headproof}

%% ---- system name -----------------------------------------------------------
\newcommand{\sys}{\mbox{OurSys}}

%% ---- measured numbers --------------------------------------------------------
%% Every measured number is \num{key}. templates/backfill/backfill.py writes the values into numbers_<lang>.tex
%% from the result files, so a rerun refills both papers. A key without a value prints in red.
\newcommand{\pending}[1]{\textcolor{red}{[\detokenize{#1}]}}
\newcommand{\setnum}[2]{\expandafter\gdef\csname rbnum@#1\endcsname{#2}}
\newcommand{\num}[1]{\ifcsname rbnum@#1\endcsname\csname rbnum@#1\endcsname\else\pending{#1}\fi}

%% ---- notation (docs/CONCEPTS.md, one concept one name one symbol) ----------
\newcommand{\budget}{B}
\newcommand{\cand}{\mathcal{C}}
\newcommand{\policy}{\pi}
